\section{Introduction}

The promise that quantum computers will solve the central problem of electronic
structure---computing the ground-state energy of a many-electron Hamiltonian, whose
configuration space grows combinatorially and whose exact solution is QMA-complete in the worst
case~\cite{schuchverstraete2009}---has driven a decade of algorithmic development~\cite{reiher2017,cerezo2021,tilly2022}. The
variational quantum eigensolver (VQE)~\cite{peruzzo2014,kandala2017} dominated the noisy-device era
but has, for the purpose of demonstrating a scaling advantage, largely run its course: barren
plateaus render deep expressive circuits untrainable~\cite{mcclean2018}, and the very structural
conditions that let a circuit avoid them tend also to render it classically
simulable~\cite{cerezo2025bp}. Sample-based quantum diagonalization
(SQD)~\cite{robledomoreno2025}, equivalently quantum-selected configuration interaction
(QSCI)~\cite{kanno2023}, took a different route: rather than optimize a circuit to \emph{be} the ground state, one uses a
shallow, noise-tolerant circuit only to \emph{sample} the electronic configurations that carry
appreciable weight, and then diagonalizes the Hamiltonian classically in the sampled determinant
subspace on high-performance hardware~\cite{robledomoreno2025,skqd2025,yoshioka2025,kirby2023}.
Within two years SQD became the de facto default of pre-fault-tolerant quantum chemistry, reaching
iron--sulfur clusters at 77 qubits with the Fugaku supercomputer in the
loop~\cite{robledomoreno2025,shirakawa2025} and extending to excited
states~\cite{extsqd2025}, adaptive state preparation~\cite{adaptqsci}, density-matrix
embedding~\cite{shajan2024}, and auxiliary-field quantum Monte Carlo~\cite{danilov2025,qcqmc,qcafqmcion2025}.

Yet the target of SQD---the ground-state \emph{energy}---is exactly where the classical frontier has
closed the gap. The broader case for a generic exponential quantum advantage in ground-state
chemistry was already found wanting: across sampled chemical space, heuristic classical methods
remain competitive and the assumed advantage rests on efficient state
preparation~\cite{leechan2023}. For SQD specifically, an extended matchgate simulator reproduces the
unitary-cluster-Jastrow--SQD workflow~\cite{extraferm}, single-layer unitary cluster Jastrow
energies are computable in polynomial time---reproducing the flagship 77-qubit experiment on a
laptop at \emph{lower} energy~\cite{belagali2026}---and classical autoregressive neural samplers
reconstruct the determinant-selection loop on a GPU~\cite{hinqs2026,arnnsci2026}. A direct
quantum-chemistry benchmark finds classical heat-bath configuration interaction more compact and
cheaper than QSCI at accessible sizes~\cite{reinholdt2025}, and a complementary lattice-model study
reports the intrinsic exponential configuration-space scaling that drives this~\cite{gaberle2026}---a
conclusion we reached independently in a companion review~\cite{review}. This mirrors the fate of the first ``quantum utility''
demonstration~\cite{kim2023}, reproduced classically within months by tensor-network~\cite{tindall2024}
and sparse-Pauli~\cite{begusic2024} methods. The lesson is quantitative: under fault-tolerant
overheads even polynomial speedups do not deliver practical advantage~\cite{babbush2021}, so a
\emph{defensible} advantage must live where the classical cost genuinely
explodes~\cite{mazzola2024}.

We argue that the \emph{defensible} quantum advantage of a sample-based method, if one exists, is not in
the energy but in the \emph{sampling} of a hard distribution, realized most cleanly for
\emph{dynamical} observables---the regime our primitive targets, a necessary precondition for any such
advantage rather than a demonstration of one. Real-time dynamics is where classical tensor networks hit the
entanglement barrier~\cite{calabresecardy2005}, where the fermion sign problem becomes a phase
problem~\cite{troyerwiese2005}, and where every
surviving beyond-classical hardware demonstration now lives---from the random-circuit-sampling
programme~\cite{arute2019} and verifiable-advantage correlators~\cite{google_echoes} to real-time
simulations of the Fermi--Hubbard model that already outrun exact classical methods on
superconducting~\cite{googlefh2025} and trapped-ion~\cite{quantinuumfh2025} processors, and to
precision estimation of prethermal dynamics beyond classical
reach~\cite{qedma2026,algorithmiq2026}. The natural dynamical targets are the
\emph{spectral functions} of correlated matter: the single-particle spectral function $A(\omega)$
measured by angle-resolved photoemission~\cite{damascelli2003}, its momentum-resolved form
$A(k,\omega)$, and the dynamical structure factor $S(q,\omega)$ probed by neutron and X-ray
scattering~\cite{squires2012,ament2011}---the central outputs of dynamical mean-field theory~\cite{dmft1996}. Classically these
are obtained by dynamical density-matrix renormalization~\cite{jeckelmann2002,nocera2016},
time-dependent DMRG~\cite{whitefeiguin2004}, and kernel-polynomial expansions~\cite{kpm2006} in one
dimension, where the density-matrix renormalization group is near-optimal~\cite{whitedmrg1992,schollwock2011}, but
they remain hard beyond it.

Crucially, no published method reconstructs a frequency-resolved single-particle spectral function from
a \emph{bitstring-sampled} quantum subspace in the charged $(N{\pm}1)$ sectors, with the Green's function
assembled classically in the Lehmann representation. The SQD/SKQD/TE-QSCI
lineage stops at energies and, at most, a handful of discrete excited-state
levels~\cite{extsqd2025,teqsci2024}; quantum Green's-function methods use Hadamard
tests~\cite{bauer2016,endo2020} or Krylov and quantum-subspace recursions
built from measured overlaps~\cite{mcclean2017qse,rungger2020,cortesgray2022,parrish2019,onqse2024,umeano2025}---the
last computing a dynamical structure factor $S(q,\omega)$ for a Kitaev spin liquid from a measured-overlap
subspace rather than sampled configurations---or hybrid quantum--classical response-and-structure
frameworks~\cite{du2026}, rather than sampled subspaces; and
the hardware real-time-observable results come from a separate
error-mitigation-plus-Trotter line~\cite{qedma2026,algorithmiq2026,quantinuum_dsf2026}. The two
closest results---a QSCI quasiparticle band structure that returns discrete
poles~\cite{ohgoe2025} and a molecular spectrum obtained from a neutral-sector autocorrelation built from
short-time quantum-evolved samples with classically projected long-time dynamics~\cite{santini2025}---return
static levels or a neutral-sector response, not the charged-sector $(N{\pm}1)$ single-particle Green's
function that our Lehmann construction assembles from the sampled subspace. We occupy that gap.

To know \emph{where} such a method is hard---and, just as important, where it is not---one must know
which resource controls its cost. Free-fermion (matchgate) circuits are classically
efficient~\cite{valiant2002,terhaldivincenzo2002,bravyi2005flo,jozsamiyake2008}; the resource that
leaves this manifold is fermionic non-Gaussianity, and every pure non-Gaussian fermionic state is a
\emph{magic} state for matchgate computation~\cite{hebenstreit2019}---the fermionic counterpart of the
non-stabilizerness that powers universal qubit computation~\cite{bravyikitaev2005,leone2022sre}. The
efficiently computable fermionic AntiFlatness (FAF)~\cite{sierant2026}, built from the Majorana
covariance matrix, quantifies that distance from the Gaussian manifold, and it is natural to hope that
this single measured number decides when the quantum sampler is essential. \emph{It does not.} The
cost of a sample-based method is its determinant-subspace size $|\mathcal S|$ (defined below); we prove
one-body fermionic magic is decoupled from it. Understanding this
decoupling---and locating the genuine hardness in the higher-order sector, where the paired
``magic-input'' regime is classically tractable for additive-error \emph{estimation} under
free-fermionic dynamics~\cite{oh2026}, though computational-basis \emph{sampling} of the same family in
a rotated frame may remain hard~\cite{fermionsampling2022}---is the second contribution of this work.

This work makes \emph{two} central contributions to quantum computation, each closing a gap that a
comprehensive survey of the literature confirms is open, and each supported by a rigorous account of
its scope.

\textbf{First, a new use of the sampled-subspace toolkit: frequency-resolved dynamical response
from bitstring-sampled subspaces---charged-sector single-particle spectra and neutral-sector structure factors.} We compute the single-particle spectral function $A(\omega)$ and, from the same
charged-$(N{\pm}1)$ sampled configurations, its momentum-resolved form $A(k,\omega)$; the density and
spin responses $S(q,\omega)$ and $S^{zz}(q,\omega)$ follow from number-conserving seeds in the neutral
$N$ sector---all purely from computational-basis bitstrings, with no Hadamard tests, no controlled unitaries, no
overlap estimation. This lifts the entire SQD/QSCI toolkit from the static energies and discrete
levels it has been confined to~\cite{robledomoreno2025,kanno2023,extsqd2025,ohgoe2025} to
continuous, frequency-resolved dynamics, and we validate it against exact diagonalization on Hubbard
chains, across nineteen molecules spanning every bonding regime, and on the IBM Heron
processor.

\textbf{Second, the resource that sets the cost---and the one that does not.} The cost of a
sample-based method is the size $|\mathcal S|$ of the determinant subspace the sampler must populate to
reach a target accuracy, and we identify what controls it by first ruling out what does not. Restricted
to the number-conserving states of chemistry, the fermionic AntiFlatness~\cite{sierant2026} collapses
to a single one-particle reduced-density-matrix invariant,
$\mathcal F_1=4\,\mathrm{tr}[\gamma(1-\gamma)]$ (equal to $2N_{\mathrm u}$ for spin-restricted natural
orbitals)---the effectively-unpaired-electron
index~\cite{takatsuka1978,staroverov2000} and the one-body-entanglement linear
entropy~\cite{gigena2020} (\S\ref{sec:resource}). Being a function of the orbital-rotation-invariant
occupation spectrum it is a Gaussian invariant, whereas $|\mathcal S|$ is basis dependent; a quantity
unchanged by the orbital rotations that sweep $|\mathcal S|$ over its whole range cannot predict it.
Across the same nineteen FCI-verified molecules one-body magic does track static, multireference
character---switching on for covalent bond-breaking, already large for inherently multireference
$\mathrm{C_2}$ and $\mathrm{O_2}$, near-Gaussian for the ionic bond of LiF, refining to fermionic
non-Gaussianity the qubit non-stabilizerness peak of molecular bonding~\cite{sarkis2025}---yet this
correlation is not even sign stable: a free-fermion determinant has $\mathcal F_1=0$ with exponentially
large $|\mathcal S|$, and increasing on-site repulsion raises $\mathcal F_1$ while \emph{lowering}
$|\mathcal S|$. One-body fermionic magic is a faithful multireference diagnostic but an unreliable
predictor of the sampling cost, which is set instead by the basis-dependent entanglement---the minimal
bond dimension $\chi$, which tracks $|\mathcal S|$ across the suite. Any genuine quantum advantage lives
one stratum above, in the generic, high-rank non-Gaussianity of the connected higher-body cumulants in
the fixed sampling basis---not the orbital-invariant higher-order antiflatness $\mathcal F_{k\ge2}$, which
we exhibit large on a classically trivial paired family---beyond the stratum whose free-fermionic
estimation primitives remain classically tractable~\cite{oh2026}.
To our knowledge this is the first time a fermionic-magic measure is connected to sample-based
diagonalization.

These contributions are delivered with the rigor the dequantization debate now
demands~\cite{leechan2023,reinholdt2025,belagali2026}: we scope the reach of the method (it
is asymptotic, and lives in two-dimensional and chemical settings, not in one dimension where tensor
networks remain near-optimal~\cite{whitedmrg1992,schollwock2011}), prove a moment-exactness and polynomial-shot bound
for the sampled reconstruction, and, on the artificial-intelligence side of the hybrid loop, show that
self-consistent configuration recovery~\cite{robledomoreno2025} improves the sampled subspace under
realistic device noise~\cite{wray2025,vaquero2026}---through valid-shot economy rather than quantum advantage---while a \emph{learned}
generative tail-completer does \emph{not} beat that classical baseline (Fig.~\ref{fig:gflow}), delimiting
the quantum--AI frontier consistent with, not counter to, the classical-simulability frontier we and
others have charted~\cite{review,reinholdt2025}. The
resulting resource diagnostics, computed from bitstrings already in hand, map where that frontier lies
for the spectral and dynamical workloads of quantum-centric
computation~\cite{robledomoreno2025,shirakawa2025}.

Beyond the specific algorithm, the reframing carries weight for both the science and the industry of
quantum computing. \emph{For the field}, it forges a first link between the resource theory of
fermionic non-Gaussianity~\cite{hebenstreit2019,sierant2026} and sample-based electronic-structure
computation~\cite{robledomoreno2025,kanno2023}, and replaces a natural but mistaken hope---that a
one-body magic number certifies where sampling is hard---with a precise statement of which resource
sets the cost and which does not. \emph{For practice}, the observables it delivers---the single-particle
spectral function, its momentum resolution, and the dynamical structure factor of correlated
matter~\cite{damascelli2003,dmft1996}---are precisely those that materials discovery, catalysis, and
correlated-electron device design demand and that the energy-only output of today's sample-based
hardware does not provide; and because the resource diagnostics are measured from bitstrings already in
hand, they chart, before further runtime is spent, where the classical frontier lies for a given
spectral or dynamical workload on the quantum-centric supercomputers now in
operation. We develop both threads in Sec.~\ref{sec:honest}.
